% =============================================================================
% Chapter 8: Model Training — ML Systems Lecture Slides
% =============================================================================
% This deck covers the training systems chapter: Iron Law of Training,
% memory breakdown, pipeline architecture, optimizations (mixed precision,
% checkpointing, FlashAttention, gradient accumulation), and scaling.
%
% IMAGES (SVGs converted to PDF):
%   - ch08-iron-law-training.pdf      - ch08-training-pipeline.pdf
%   - ch08-training-memory-breakdown.pdf
%   - ch08-mixed-precision-flow.pdf   - ch08-gradient-checkpointing.pdf
%   - ch08-optimization-walkthrough.pdf
%   - ch08-data-parallelism.pdf
% =============================================================================
\documentclass[aspectratio=169, 12pt]{beamer}
\usepackage{../../assets/beamerthememlsys}

\mlsyssetup{
  volume   = {Volume I},
  chapter  = {Chapter 8},
  logo = {../../assets/img/logo-mlsysbook.png},
  instlogo = {../../assets/img/logo-harvard.png},
  chaptertitle = {Model Training},
}

% --- Fonts ---
\usepackage{fontspec}
\setsansfont{Helvetica Neue}[
  BoldFont={Helvetica Neue Bold},
  ItalicFont={Helvetica Neue Italic},
  BoldItalicFont={Helvetica Neue Bold Italic},
]
\setmonofont{JetBrains Mono}[Scale=0.85]

% --- Packages ---
\usepackage{booktabs}
\usepackage{amsmath}

% --- Image paths ---
\graphicspath{
  {images/}
}

% --- Chapter-specific macros ---
\newcommand{\DAM}{D\raisebox{0.08em}{\tiny$\bullet$}A\raisebox{0.08em}{\tiny$\bullet$}M}

% --- Helper: safe image include ---
\newcommand{\safeimg}[2][width=\textwidth,keepaspectratio]{%
  \IfFileExists{images/#2}{\includegraphics[#1]{#2}}{%
    \IfFileExists{#2}{\includegraphics[#1]{#2}}{%
      \fbox{\parbox[c][2.5cm][c]{0.85\linewidth}{\centering\footnotesize\textcolor{midgray}{[Missing image]}}}%
    }%
  }%
}

% --- Section count for navigation (must match actual \section{} count) ---
\setcounter{mlsystotalsections}{7}

\title{Model Training}
\author{Vijay Janapa Reddi}
\institute{Harvard University}
\date{}

\begin{document}

% =============================================================================
% TITLE SLIDE
% =============================================================================
\mlsystitle{Model Training}{The Physics of Learning at Scale}{cover_ai_training.png}

% =============================================================================
% LEARNING OBJECTIVES
% =============================================================================
\begin{frame}{Learning Objectives}
\note{[2 min] Read through objectives. Emphasize that this chapter bridges
framework internals (Ch7) with the physical reality of training at scale.
Ask: ``How many of you have waited hours for a model to train and wondered
why it takes so long?''}

\footnotesize
\begin{enumerate}\setlength\itemsep{0pt}
  \item Explain the \textbf{Iron Law of Training Performance} and identify which term each optimization targets
  \item Calculate \textbf{FLOPs} and \textbf{memory footprints} (activations, optimizer states) for training
  \item Compare \textbf{SGD}, \textbf{Adam}, and \textbf{AdamW} by convergence speed, memory, and cost
  \item Apply \textbf{arithmetic intensity analysis} to classify ops as compute-bound or memory-bound
  \item Apply \textbf{mixed-precision training}, \textbf{checkpointing}, and \textbf{gradient accumulation}
  \item Construct efficient pipelines using \textbf{data prefetching} and \textbf{pipeline overlapping}
  \item Analyze when single-machine training becomes infeasible and evaluate \textbf{data parallelism}
\end{enumerate}

\end{frame}

% =============================================================================
\section{Training Costs}
% =============================================================================

\begin{frame}{Why Training Costs Millions}
\note{[3 min] Open with the cost asymmetry. GPT-2 inference: fraction of a cent.
GPT-2 training: \$50K. GPT-4 training: \$100M estimated. This million-to-one
gap governs who can participate in AI. Ask: ``Who can afford to train a
frontier model?''}

\small
\begin{columns}[T]
  \begin{column}{0.55\textwidth}
    Running a model costs \textbf{pennies}.\\
    Training a model costs \textbf{millions}.

    \vspace{0.2cm}
    \renewcommand{\arraystretch}{1.2}
    {\footnotesize
    \begin{tabular}{@{}lrr@{}}
      \toprule
      \textbf{Model} & \textbf{Inference} & \textbf{Training} \\
      \midrule
      GPT-2  & $<$\$0.01 & \$50K (2019) \\
      GPT-3  & \$0.02/call & \$4.6M (2020) \\
      GPT-4  & \$0.03/call & $\approx$\$100M (2023) \\
      \bottomrule
    \end{tabular}
    }

    \vspace{0.15cm}
    \begin{mlsyscard}{crimson}
    {\footnotesize \textbf{Million-to-one cost asymmetry:} Training requires billions of forward passes, each paired with a backward pass that costs $\approx$2$\times$ the forward pass.}
    \end{mlsyscard}
  \end{column}
  \begin{column}{0.42\textwidth}
    \vspace{0.2cm}
    {\footnotesize
    \textbf{Three scaling challenges:}
    \begin{itemize}\setlength\itemsep{0pt}
      \item \textcolor{computestroke}{\textbf{Computational intensity}}: Sustained petaFLOPS for days
      \item \textcolor{datastroke}{\textbf{Memory pressure}}: Weights + gradients + optimizer + activations
      \item \textcolor{errorstroke}{\textbf{Data dependencies}}: Sequential gradient updates
    \end{itemize}
    }
  \end{column}
\end{columns}

\end{frame}

\begin{frame}{The 3AM Gradient Explosion}
\note{[2 min] Tell this as a story. A team trains a 7B model for 4 days.
At 3AM, loss spikes to NaN. They lost 4 days and \$5000. The cause: FP16
overflow from an outlier batch. The fix: checkpointing, gradient clipping,
BF16. This motivates the rest of the chapter.}

\scriptsize
\begin{columns}[T]
  \begin{column}{0.55\textwidth}
    \textbf{Scenario:} Training a 7B LLM. Loss smooth for 4 days.\\
    \textbf{3:00 AM:} Loss spikes to \textbf{NaN}. Crash. Lost 4 days (\$5K).

    \vspace{0.05cm}
    \textbf{Physics:} Outlier batch exceeded FP16 range ($>$65,504).\\
    Gradients $\to$ Inf $\to$ all weights NaN.

    \vspace{0.05cm}
    \textbf{Systems Fix:}
    \begin{itemize}\setlength\itemsep{0pt}
      \item \textbf{Checkpointing}: Save hourly (lose 1hr, not 4 days)
      \item \textbf{Gradient clipping}: Cap gradient norms
      \item \textbf{BF16}: FP32 range in 16 bits; overflows rare
    \end{itemize}
  \end{column}
  \begin{column}{0.42\textwidth}
    \renewcommand{\arraystretch}{1.1}
    \begin{tabular}{@{}ll@{}}
      \toprule
      \textbf{Format} & \textbf{Max Value} \\
      \midrule
      FP32 & $3.4 \times 10^{38}$ \\
      FP16 & 65,504 \\
      BF16 & $3.4 \times 10^{38}$ \\
      \bottomrule
    \end{tabular}

    \vspace{0.1cm}
    \mlsysalert{Key:}{BF16 trades precision for range---FP32 exponent in 16 bits.}
  \end{column}
\end{columns}

\end{frame}

% =============================================================================
\section{Iron Law}
% =============================================================================

\mlsysfocus{The Iron Law of Training Performance}{%
$T_{\text{train}} = \dfrac{O}{R_{\text{peak}} \times \eta}$%
}

\begin{frame}{Iron Law: Three Levers}
\note{[3 min] Walk through each term. O = total FLOPs. R\_peak = hardware
throughput (set by procurement and precision). eta = utilization (the
engineering target). GPT-3 achieved eta = 0.45. Current target: >0.55.
Ask: ``Which lever is cheapest to improve?''}

% --- Layout: FULL-WIDTH diagram ---
\centering
\safeimg[width=0.92\textwidth,height=4.8cm,keepaspectratio]{ch08-iron-law-training.pdf}

\end{frame}

\begin{frame}{Iron Law: Optimization Mapping}
\note{[2 min] Every technique in this chapter maps to one of the three terms.
Students should be able to classify any optimization by which term it targets.
This table is the organizing framework for the rest of the deck.}

\scriptsize
\renewcommand{\arraystretch}{1.1}
\begin{tabular}{@{}llp{5.5cm}@{}}
  \toprule
  \textbf{Technique} & \textbf{Term} & \textbf{Mechanism} \\
  \midrule
  \textbf{Mixed Precision} & $R_{\text{peak}}$ $\uparrow$ & Tensor Cores at 16$\times$ higher FLOP/s \\
  \textbf{Data Prefetching} & $\eta$ $\uparrow$ & Reduces idle time waiting for data \\
  \textbf{Grad Checkpointing} & $O$ $\uparrow$ & Adds recomputation, enables larger models \\
  \textbf{Grad Accumulation} & $\eta$ $\uparrow$ & Maintains batch parallelism efficiency \\
  \textbf{Operator Fusion} & $\eta$ $\uparrow$ & Reduces memory bandwidth bottlenecks \\
  \textbf{FlashAttention} & $O$ $\downarrow$, $\eta$ $\uparrow$ & IO-aware tiling reduces FLOPs, improves access \\
  \bottomrule
\end{tabular}

\vspace{0.1cm}
{\footnotesize\mlsysinsight{Diagnostic discipline:}{Identify which term dominates \emph{before} choosing an optimization.}}

\end{frame}

% --- ACTIVE LEARNING 1: Predict ---
\begin{frame}{Predict: What Dominates Training Memory?}
\note{[2 min] Prediction exercise. Give students 60 seconds. Most will guess
model weights. The surprising answer: optimizer state and activations
dominate. This primes the memory breakdown slide.}

\centering
\vspace{1.0cm}
{\Large\bfseries Think--Write--Share}

\vspace{0.6cm}
{\large A 7B-parameter model in FP16 weighs 14 GB.\\[0.2cm]
How much total GPU memory does \textbf{training} require?}

\vspace{0.5cm}
{\normalsize Write your estimate. \textcolor{midgray}{(60 seconds)}}

\vspace{0.5cm}
{\small\textcolor{lightgray}{Hint: Think about what else must live in GPU memory besides weights.}}

\end{frame}

% =============================================================================
\section{Memory \& Math}
% =============================================================================

\begin{frame}{Training vs.\ Inference: The 4$\times$ Multiplier}
\note{[3 min] Reveal after prediction. The 4x multiplier is the most
important number: weights (14 GB) + gradients (14 GB) + Adam m (14 GB) +
Adam v (14 GB) = 56 GB minimum. Plus activations. GPT-2 XL in FP32:
89 GB total. V100 has 32 GB. This is why optimization matters.
Common error: students think weights dominate memory.}

% --- Layout: FULL-WIDTH diagram ---
\centering
\safeimg[width=0.88\textwidth,height=4.8cm,keepaspectratio]{ch08-training-memory-breakdown.pdf}

\vspace{0.1cm}
\mlsysalert{Key:}{A 7B model that \emph{runs} inference on one GPU requires \emph{multiple} GPUs for training.}

\end{frame}

\begin{frame}{Optimizer Memory: SGD vs.\ Adam}
\note{[2 min] SGD stores only gradients (1 extra copy). Adam stores m and v
(3 extra copies). Adam converges in ~1/3 the iterations but costs 3x memory.
This is a binding constraint for large models.
Ask: ``When would you choose SGD over Adam?''}

\small
\begin{columns}[T]
  \begin{column}{0.55\textwidth}
    \renewcommand{\arraystretch}{1.2}
    {\footnotesize
    \begin{tabular}{@{}lrrr@{}}
      \toprule
      & \textbf{SGD} & \textbf{SGD+M} & \textbf{Adam} \\
      \midrule
      Extra state/param & 0 & 1 & 2 \\
      Memory multiplier & 2$\times$ & 3$\times$ & 4$\times$ \\
      7B model (FP16+FP32) & 28 GB & 42 GB & 56 GB \\
      Convergence speed & Slow & Medium & Fast \\
      \bottomrule
    \end{tabular}
    }

    \vspace{0.2cm}
    \begin{mlsyscard}{routingstroke}
    {\footnotesize \textbf{Adam} converges in $\sim$1/3 the iterations of SGD but requires \textbf{3$\times$} the memory for per-parameter state.}
    \end{mlsyscard}
  \end{column}
  \begin{column}{0.42\textwidth}
    \vspace{0.1cm}
    {\footnotesize
    \textbf{Adam update rule:}\\[0.1cm]
    $m_t = \beta_1 m_{t-1} + (1-\beta_1) g_t$\\
    $v_t = \beta_2 v_{t-1} + (1-\beta_2) g_t^2$\\
    $\theta_t = \theta_{t-1} - \alpha \frac{\hat{m}_t}{\sqrt{\hat{v}_t} + \epsilon}$\\[0.2cm]
    Two state vectors ($m$, $v$) per parameter.\\[0.1cm]
    \textbf{AdamW} adds weight decay decoupled from the gradient, improving generalization.
    }
  \end{column}
\end{columns}

\end{frame}

\begin{frame}{Activation Memory: The Hidden Cost}
\note{[2 min] Activations scale with batch size and depth. GPT-2: 65 GB
of activations in FP32 (seq\_len=1024, batch=32, 48 layers). This often
exceeds optimizer state. Activation checkpointing trades compute for memory
(covered later). If short: just show the formula and GPT-2 number.}

\small
\begin{columns}[T]
  \begin{column}{0.55\textwidth}
    Activations must be stored from the forward pass for use in the backward pass.

    \vspace{0.15cm}
    {\footnotesize
    $$\text{Act. Memory} \propto B \times S \times H \times L$$

    \begin{tabular}{@{}ll@{}}
      $B$ & batch size \\
      $S$ & sequence length \\
      $H$ & hidden dimension \\
      $L$ & number of layers \\
    \end{tabular}
    }

    \vspace{0.1cm}
    \mlsysalert{GPT-2 XL:}{48 layers, batch=32: activations $\approx$ \textbf{65 GB}.}
  \end{column}
  \begin{column}{0.42\textwidth}
    {\footnotesize
    \textbf{GPT-2 FP32 breakdown:}

    \vspace{0.05cm}
    \renewcommand{\arraystretch}{1.05}
    \begin{tabular}{@{}lr@{}}
      \toprule
      \textbf{Component} & \textbf{Size} \\
      \midrule
      Weights & 6 GB \\
      Gradients & 6 GB \\
      Adam states & 12 GB \\
      Activations & 65 GB \\
      \midrule
      \textbf{Total} & \textbf{89 GB} \\
      \bottomrule
    \end{tabular}

    \vspace{0.1cm}
    Activations = \textbf{73\%} of total memory.
    }
  \end{column}
\end{columns}

\end{frame}

\begin{frame}{Arithmetic Intensity: Compute-Bound vs.\ Memory-Bound}
\note{[2 min] Arithmetic intensity = FLOPs per byte loaded. If intensity
is below the ridge point, the operation is memory-bound. Most training ops
(ReLU, LayerNorm, Softmax) are memory-bound. Only large GEMMs are
compute-bound. Increasing batch size moves operations rightward.
This is the roofline model from Ch1 applied to training.}

\small
\begin{columns}[T]
  \begin{column}{0.55\textwidth}
    $$\text{Arithmetic Intensity} = \frac{\text{FLOPs}}{\text{Bytes loaded}}$$

    \vspace{0.1cm}
    {\footnotesize
    \renewcommand{\arraystretch}{1.15}
    \begin{tabular}{@{}lrl@{}}
      \toprule
      \textbf{Operation} & \textbf{AI} & \textbf{Regime} \\
      \midrule
      ReLU & $<$1 & \textcolor{errorstroke}{Memory-bound} \\
      LayerNorm & $\sim$5 & \textcolor{errorstroke}{Memory-bound} \\
      Softmax & $\sim$8 & \textcolor{errorstroke}{Memory-bound} \\
      GEMM (batch=1) & $\sim$20 & \textcolor{routingstroke}{Transitional} \\
      GEMM (batch=256) & $>$200 & \textcolor{datastroke}{Compute-bound} \\
      \bottomrule
    \end{tabular}
    }
  \end{column}
  \begin{column}{0.42\textwidth}
    \vspace{0.2cm}
    \begin{mlsyscard}{computestroke}
    {\footnotesize \textbf{Most training ops are memory-bound} except large-batch GEMMs.\\[0.1cm]
    Increasing batch size moves operations toward compute-bound.}
    \end{mlsyscard}

    \vspace{0.15cm}
    {\footnotesize V100 ridge point: 139 FLOPs/byte\\
    \textcolor{midgray}{Below ridge $\to$ memory-bound}}
  \end{column}
\end{columns}

\end{frame}

% =============================================================================
\section{Pipeline}
% =============================================================================

\begin{frame}{The Staged Training Pipeline}
\note{[3 min] Training is a 4-stage pipeline. Data loading (CPU), PCIe
transfer, forward/backward (GPU), gradient sync (NVLink). Any mismatch
creates ``accelerator bubbles'' --- idle GPU time. The goal: overlap stages.
Walk through the naive vs pipelined diagram.}

% --- Layout: FULL-WIDTH diagram ---
\centering
\safeimg[width=0.92\textwidth,height=4.8cm,keepaspectratio]{ch08-training-pipeline.pdf}

\vspace{0.1cm}
\mlsysinsight{Engineering goal:}{Overlap all stages so the GPU never waits for data.}

\end{frame}

\begin{frame}{Four Pipeline Stages}
\note{[2 min] Detail each stage. Data loading: CPU workers decode from NVMe.
Transfer: PCIe DMA to GPU memory. Compute: forward + backward on GPU.
Sync: gradient averaging via NVLink. The slowest stage determines throughput.
Common error: optimizing GPU compute when data loading is the bottleneck.}

\footnotesize
\renewcommand{\arraystretch}{1.15}
\begin{tabular}{@{}llll@{}}
  \toprule
  \textbf{Stage} & \textbf{Hardware} & \textbf{Bottleneck} & \textbf{Mitigation} \\
  \midrule
  \textbf{1. Data Loading} & CPU / NVMe & Decode + augment & Multi-worker prefetch \\
  \textbf{2. Host$\to$Device} & PCIe bus & Transfer bandwidth & Pin memory, async DMA \\
  \textbf{3. Fwd/Bwd Pass} & GPU compute & Matrix ops & Mixed precision, fusion \\
  \textbf{4. Grad Sync} & NVLink & Communication & Gradient overlap \\
  \bottomrule
\end{tabular}

\vspace{0.2cm}
\begin{columns}[T]
  \begin{column}{0.48\textwidth}
    \begin{mlsyscard}{errorstroke}
    {\footnotesize \textbf{Accelerator Bubbles}\\
    GPU idle while waiting for data = 0\% utilization. The ``Machine'' axis waits for ``Data.''}
    \end{mlsyscard}
  \end{column}
  \begin{column}{0.48\textwidth}
    \begin{mlsyscard}{datastroke}
    {\footnotesize \textbf{Prefetching}\\
    Load batch $N+1$ while GPU processes batch $N$. Requires async data loaders.}
    \end{mlsyscard}
  \end{column}
\end{columns}

\end{frame}

% --- ACTIVE LEARNING 2: Exercise ---
\begin{frame}{Your Turn: Memory Budget}
\note{[3 min] Give students 90 seconds. Answer: Weights=2 GB + Grads=2 GB +
Adam=4 GB + Activations (B*S*H*L*4 bytes). At batch 16, activations are
roughly 6 GB. Total ~14 GB. Fits on V100 (32 GB). At batch 64, activations
quadruple to ~24 GB. Total ~32 GB. Barely fits.
Ask: ``What happens if we increase batch size to 128?''}

\small
\begin{columns}[T]
  \begin{column}{0.55\textwidth}
    {\normalsize\bfseries Calculate Training Memory}

    \vspace{0.2cm}
    A model has \textbf{1B parameters} trained with Adam in FP16+FP32:
    \begin{itemize}\setlength\itemsep{0pt}
      \item Weights: 1B $\times$ 2 bytes = ?
      \item Gradients: 1B $\times$ 2 bytes = ?
      \item Adam states: 1B $\times$ 4 bytes = ?
      \item Activations at batch 16 $\approx$ 6 GB
    \end{itemize}

    \vspace{0.1cm}
    \textbf{Calculate} total memory. Does it fit on a V100 (32 GB)?

    {\footnotesize\textcolor{midgray}{(90 seconds --- then compare with a neighbor)}}
  \end{column}
  \begin{column}{0.42\textwidth}
    \pause
    \begin{mlsyscard}{datastroke}
    \textbf{Solution:}\\[0.1cm]
    {\footnotesize
    Weights: 2 GB\\
    Gradients: 2 GB\\
    Adam (m+v): 4 GB\\
    Activations: $\sim$6 GB\\[0.1cm]
    Total $\approx$ \textbf{14 GB}\\[0.1cm]
    \textcolor{datastroke}{\textbf{Fits on V100!}}\\[0.1cm]
    At batch 64: activations $\approx$24 GB\\
    Total $\approx$ 32 GB --- barely fits.
    }
    \end{mlsyscard}
  \end{column}
\end{columns}

\end{frame}

% =============================================================================
\section{Optimizations}
% =============================================================================

\begin{frame}{Mixed-Precision Training}
\note{[3 min] FP16 computation with FP32 accumulation. Walk through the flow:
FP32 master weights -> cast to FP16 -> forward -> loss scale -> backward ->
unscale -> FP32 accumulate. Loss scaling prevents gradient underflow in FP16.
BF16 eliminates the need for loss scaling. Result: ~2x throughput, ~2x memory.}

% --- Layout: FULL-WIDTH diagram ---
\centering
\safeimg[width=0.92\textwidth,height=4.8cm,keepaspectratio]{ch08-mixed-precision-flow.pdf}

\end{frame}

\begin{frame}{Mixed Precision: The Numbers}
\note{[2 min] Quantitative impact. FP16 on V100 Tensor Cores: 125 TFLOPS
vs 15.7 TFLOPS FP32 = nearly 8x throughput. Memory halved. BF16 simplifies
the pipeline by matching FP32's dynamic range.
Common error: enabling mixed precision without loss scaling.}

\small
\begin{columns}[T]
  \begin{column}{0.55\textwidth}
    \renewcommand{\arraystretch}{1.2}
    {\footnotesize
    \begin{tabular}{@{}lrrr@{}}
      \toprule
      \textbf{Metric} & \textbf{FP32} & \textbf{FP16} & \textbf{BF16} \\
      \midrule
      V100 TFLOPS & 15.7 & 125 & --- \\
      Bytes/param & 4 & 2 & 2 \\
      Dynamic range & $10^{38}$ & 65K & $10^{38}$ \\
      Loss scaling & No & Yes & No \\
      Accuracy impact & --- & $<$1\% & $<$1\% \\
      \bottomrule
    \end{tabular}
    }

    \vspace{0.15cm}
    \mlsysconcept{Iron Law target:}{$R_{\text{peak}}$ --- Tensor Cores operate at up to 16$\times$ higher throughput.}
  \end{column}
  \begin{column}{0.42\textwidth}
    \vspace{0.1cm}
    \begin{mlsyscard}{computestroke}
    {\footnotesize \textbf{Impact:}\\[0.05cm]
    $\sim$2$\times$ throughput\\
    $\sim$2$\times$ memory savings\\
    $<$1\% accuracy loss\\[0.1cm]
    BF16 = ``best of both worlds'':\\
    FP32 range + FP16 speed}
    \end{mlsyscard}
  \end{column}
\end{columns}
\mlsyscite{Micikevicius et al., ``Mixed Precision Training,'' ICLR 2018}

\end{frame}

\begin{frame}{Gradient Checkpointing}
\note{[3 min] Standard backprop stores all L layers' activations. Checkpointing
stores only sqrt(L) and recomputes the rest. Trade: +33\% compute for up to
10x memory reduction. GPT-2: 48 layers -> 7 checkpoints. Walk through the
side-by-side comparison. If short: focus on the trade-off numbers.}

% --- Layout: FULL-WIDTH comparison diagram ---
\centering
\safeimg[width=0.92\textwidth,height=4.8cm,keepaspectratio]{ch08-gradient-checkpointing.pdf}

\end{frame}

\begin{frame}{FlashAttention: IO-Aware Computation}
\note{[2 min] Standard attention materializes the full NxN matrix in HBM.
FlashAttention tiles the computation, keeping partial results in SRAM.
Never writes the NxN matrix to HBM. Result: 2-4x speedup, 4x longer
sequences. This is the principle of IO-aware algorithm design.
Iron Law: reduces O and improves eta.}

\footnotesize
\begin{columns}[T]
  \begin{column}{0.55\textwidth}
    \textbf{Standard Attention:}\\
    Materializes $n \times n$ matrix in HBM\\
    Memory: $O(n^2)$ --- 1024 tokens $\to$ 4 MB/head

    \vspace{0.1cm}
    \textbf{FlashAttention:}\\
    Tiles computation in SRAM (20 MB)\\
    Never writes full $n \times n$ to HBM

    \vspace{0.1cm}
    \begin{mlsyscard}{datastroke}
    \textbf{Impact:} 2--4$\times$ speedup, 4$\times$ longer sequences.\\
    Converts attention from \textcolor{errorstroke}{memory-bound} to \textcolor{datastroke}{compute-bound}.
    \end{mlsyscard}
  \end{column}
  \begin{column}{0.42\textwidth}
    \renewcommand{\arraystretch}{1.15}
    \begin{tabular}{@{}lll@{}}
      \toprule
      & \textbf{Standard} & \textbf{Flash} \\
      \midrule
      HBM reads & $O(n^2)$ & $O(n^2/M)$ \\
      Memory & $O(n^2)$ & $O(n)$ \\
      Speedup & 1$\times$ & 2--4$\times$ \\
      Max seq & 1K--2K & 4K--16K \\
      \bottomrule
    \end{tabular}

    \vspace{0.1cm}
    \mlsysconcept{Iron Law:}{$O$ $\downarrow$ and $\eta$ $\uparrow$ (IO-aware tiling).}
  \end{column}
\end{columns}
\mlsyscite{Dao et al., ``FlashAttention,'' NeurIPS 2022}

\end{frame}

\begin{frame}{Gradient Accumulation}
\note{[2 min] When a large batch doesn't fit in memory, split it into
micro-batches. Accumulate gradients across micro-batches before updating.
Physical batch = small (fits in memory). Effective batch = large (good
convergence). Trades utilization for memory: 5-10\% overhead from multiple
forward/backward passes per update. If short: just cover the concept.}

\small
\begin{columns}[T]
  \begin{column}{0.55\textwidth}
    \textbf{Problem:} Batch size 256 gives best convergence\\
    but requires too much activation memory.

    \vspace{0.15cm}
    \textbf{Solution:} Split into 8 micro-batches of 32:

    \vspace{0.1cm}
    {\footnotesize
    \begin{enumerate}\setlength\itemsep{0pt}
      \item Forward/backward on micro-batch (size 32)
      \item \emph{Accumulate} gradients (do not update)
      \item Repeat for all 8 micro-batches
      \item One optimizer step with averaged gradients
    \end{enumerate}
    }

    \vspace{0.1cm}
    \mlsysconcept{Iron Law target:}{$\eta$ --- maintains effective batch efficiency within memory limits.}
  \end{column}
  \begin{column}{0.42\textwidth}
    \vspace{0.1cm}
    {\footnotesize
    \renewcommand{\arraystretch}{1.15}
    \begin{tabular}{@{}lr@{}}
      \toprule
      \textbf{Parameter} & \textbf{Value} \\
      \midrule
      Effective batch & 256 \\
      Physical batch & 32 \\
      Accumulation steps & 8 \\
      Memory reduction & $\sim$4$\times$ \\
      Overhead & 5--10\% \\
      \bottomrule
    \end{tabular}
    }

    \vspace{0.15cm}
    \begin{mlsyscard}{routingstroke}
    {\footnotesize Same convergence as large batch, but \textbf{fits in memory.}}
    \end{mlsyscard}
  \end{column}
\end{columns}

\end{frame}

\begin{frame}{GPT-2 Optimization Walkthrough}
\note{[3 min] Walk through the waterfall: 89 GB baseline -> 51 GB (mixed
precision) -> 32 GB (checkpointing) -> 28 GB (gradient accumulation).
V100 has 32 GB. Each step targets a different Iron Law term. The key
lesson: no single technique is sufficient; they compose.
Ask: ``Which step had the biggest impact?''}

% --- Layout: FULL-WIDTH waterfall diagram ---
\centering
\safeimg[width=0.92\textwidth,height=4.8cm,keepaspectratio]{ch08-optimization-walkthrough.pdf}

\end{frame}

% --- ACTIVE LEARNING 3: Discussion ---
\begin{frame}{Discussion: Which Term Would You Optimize First?}
\note{[3 min] Turn-and-talk. A training job takes 10 days. Profiling shows:
GPU utilization 35\%, FP32 precision, no checkpointing.
No single right answer, but the diagnostic order matters:
1. Mixed precision (R\_peak) --- easiest to enable
2. Profiling shows data loading at 40\% idle (eta) --- add prefetching
3. Memory pressure --- add checkpointing (O tradeoff)
Cold-call 2-3 pairs.}

\centering
\vspace{0.6cm}
{\large\bfseries Turn and Talk \textcolor{midgray}{(90 seconds)}}

\vspace{0.5cm}
{\large Your training job takes \textbf{10 days}.\\
Profiling reveals: GPU utilization = \textbf{35\%},\\
FP32 precision, no checkpointing.\\[0.3cm]
\alert{In what order would you apply optimizations?}}

\vspace{0.5cm}
\begin{columns}[c]
  \begin{column}{0.22\textwidth}\centering\small Mixed\\Precision\end{column}
  \begin{column}{0.22\textwidth}\centering\small Data\\Prefetch\end{column}
  \begin{column}{0.22\textwidth}\centering\small Gradient\\Checkpoint\end{column}
  \begin{column}{0.22\textwidth}\centering\small Grad\\Accumulation\end{column}
\end{columns}

\end{frame}

% =============================================================================
\section{Scaling}
% =============================================================================

\begin{frame}{Data Parallelism: Replicate Model, Split Data}
\note{[3 min] Walk through the diagram: mini-batch splits across GPUs.
Each GPU has a full model copy, processes its shard independently, then
AllReduce averages gradients. Communication cost: ~2x model size per step.
NVLink: 4 GPUs gives ~3.6x speedup. PCIe: only ~2.5x.
Ask: ``Why doesn't 4 GPUs give exactly 4x speedup?''}

% --- Layout: FULL-WIDTH diagram ---
\centering
\safeimg[width=0.88\textwidth,height=4.5cm,keepaspectratio]{ch08-data-parallelism.pdf}

\vspace{0.1cm}
\mlsysalert{Trade-off:}{Communication overhead grows with model size. Interconnect bandwidth is the binding constraint.}

\end{frame}

\begin{frame}{The Communication Tax}
\note{[2 min] Show the scaling curves. Ideal = linear. Compute-bound
workloads (ResNet): ~95\% efficient. Balanced (LLM + NVLink): ~85\%.
Bandwidth-bound (slow network): ~60\%. The gap is the communication tax.
Key insight: always exhaust single-machine optimizations before scaling.}

\small
\begin{columns}[T]
  \begin{column}{0.55\textwidth}
    {\footnotesize
    Effective speedup with $N$ GPUs:
    $$S(N) = \frac{N}{1 + (N-1) \cdot r}$$
    where $r$ = fraction of time spent communicating.
    }

    \vspace{0.1cm}
    \renewcommand{\arraystretch}{1.15}
    {\footnotesize
    \begin{tabular}{@{}llr@{}}
      \toprule
      \textbf{Workload} & \textbf{$r$} & \textbf{4-GPU} \\
      \midrule
      Compute-bound (ResNet) & 0.05 & 3.5$\times$ \\
      Balanced (LLM+NVLink) & 0.15 & 2.8$\times$ \\
      BW-bound (slow net) & 0.40 & 1.7$\times$ \\
      \bottomrule
    \end{tabular}
    }
  \end{column}
  \begin{column}{0.42\textwidth}
    \vspace{0.1cm}
    \begin{mlsyscard}{errorstroke}
    {\footnotesize \textbf{Scaling pitfall:}\\
    Small models on 8 GPUs spend 30--50\% synchronizing gradients.\\[0.1cm]
    A well-optimized single A100 can \textbf{outperform} a poorly-configured 8-GPU cluster.}
    \end{mlsyscard}

    \vspace{0.1cm}
    \mlsysconcept{Rule:}{Exhaust single-machine optimizations before distributing.}
  \end{column}
\end{columns}

\end{frame}

\begin{frame}{When to Scale: The Physical Ceiling}
\note{[2 min] Three signals that single-machine training is exhausted:
1. Model too large (70B = 140 GB FP16 weights alone).
2. Training too slow (years on one GPU).
3. Dataset too large (petabyte scale).
Multi-GPU adds: aggregate memory, throughput, and I/O bandwidth.
But also adds: communication overhead, debugging complexity, failure handling.}

\small
\begin{columns}[T]
  \begin{column}{0.48\textwidth}
    \textbf{Three signals to scale out:}

    \vspace{0.15cm}
    {\footnotesize
    \textcolor{errorstroke}{\textbf{1. Memory exhaustion}}\\
    70B model = 140 GB weights (FP16)\\
    $>$ any single GPU (80 GB max)

    \vspace{0.1cm}
    \textcolor{routingstroke}{\textbf{2. Unacceptable duration}}\\
    GPT-2 on 1 GPU: $\sim$months\\
    On 8 GPUs: $\sim$weeks

    \vspace{0.1cm}
    \textcolor{computestroke}{\textbf{3. Dataset scale}}\\
    Petabyte datasets require\\
    distributed storage bandwidth
    }
  \end{column}
  \begin{column}{0.48\textwidth}
    \vspace{0.2cm}
    {\footnotesize
    \renewcommand{\arraystretch}{1.15}
    \begin{tabular}{@{}ll@{}}
      \toprule
      \textbf{Strategy} & \textbf{Splits} \\
      \midrule
      Data parallel & Data (replicates model) \\
      Model parallel & Layers across GPUs \\
      Pipeline parallel & Stages across GPUs \\
      Tensor parallel & Ops within a layer \\
      \bottomrule
    \end{tabular}
    }

    \vspace{0.15cm}
    \begin{mlsyscard}{crimson}
    {\footnotesize \textbf{Conservation of Complexity:}\\
    Eliminating single-machine bottlenecks introduces new communication bottlenecks.}
    \end{mlsyscard}
  \end{column}
\end{columns}

\end{frame}

% =============================================================================
\section{Wrap-Up}
% =============================================================================

\begin{frame}{Fallacies}
\note{[2 min] Three common misconceptions with quantitative evidence.}

\small
\textbf{Fallacy:} \textit{Larger models always yield better performance.}\\
{\footnotesize A 20B model trained on insufficient data ($<$500M examples) \emph{underperforms} a 7B model by 3--5\% due to overfitting. Model capacity must match dataset size.}

\vspace{0.15cm}
\textbf{Fallacy:} \textit{Hyperparameters transfer from small-scale to large-scale training.}\\
{\footnotesize Learning rate at batch 32 does not work at batch 4096. Linear scaling rule: multiply LR by the batch ratio. Ignoring this causes divergence 2--5 days into a multi-week run.}

\vspace{0.15cm}
\textbf{Fallacy:} \textit{Mixed precision is a simple toggle.}\\
{\footnotesize FP16 without loss scaling causes gradient underflow. A model training for 72 hours can diverge at step 40K due to accumulated numerical errors. Always validate on representative workloads.}

\end{frame}

\begin{frame}{Pitfalls}
\note{[2 min] Three operational pitfalls. If short: cover just the first two.}

\small
\textbf{Pitfall:} \textit{Assuming distributed training automatically accelerates development.}\\
{\footnotesize Small models on 8 GPUs spend 30--50\% of time synchronizing. A well-optimized single A100 (\textbf{48 hours}) outperforms a poorly-configured 8-GPU cluster (\textbf{36 hours} but 8$\times$ cost). Always exhaust single-machine optimizations first.}

\vspace{0.15cm}
\textbf{Pitfall:} \textit{Optimizing memory and computation independently.}\\
{\footnotesize Utilization drops from 70\% at batch 256 to 25--35\% at batch 16. Gradient accumulation trades 5--10\% efficiency for 4$\times$ memory reduction. Tuning independently extends training by 15--30\%.}

\vspace{0.15cm}
\textbf{Pitfall:} \textit{Neglecting data pipeline optimization.}\\
{\footnotesize Data loading creates 20--40\% idle time. Prefetching reduces wall-clock time by 25--35\%. Profile \emph{before} assuming the GPU is the bottleneck.}

\end{frame}

% --- RETRIEVAL PRACTICE ---
\begin{frame}{What Were the Key Ideas?}
\note{[2 min] Retrieval practice. Students write 90 seconds, no notes.
Walk around the room. Do NOT show next slide yet.}

\centering
\vspace{1.5cm}
{\Large\bfseries Close your notes.}

\vspace{0.8cm}
{\large Write down the \textbf{4 most important concepts} from today.}

\vspace{0.8cm}
{\normalsize\textcolor{midgray}{90 seconds --- no peeking.}}

\end{frame}

\begin{frame}{Key Takeaways}
\note{[2 min] Reveal. Walk through each bullet. Emphasize the Iron Law
as the diagnostic framework. The 4x memory multiplier. Profile before
optimize. Compose optimizations systematically.}

\scriptsize
\begin{itemize}\setlength\itemsep{0pt}
  \item \textbf{Iron Law}: $T_{\text{train}} = O / (R_{\text{peak}} \times \eta)$. Every optimization targets one of these three terms.
  \item \textbf{4$\times$ memory multiplier}: Adam adds 3 extra copies per parameter. Activations scale with batch $\times$ depth.
  \item \textbf{Mixed precision}: 2$\times$ throughput, 2$\times$ memory savings, $<$1\% accuracy impact. BF16 eliminates loss scaling.
  \item \textbf{FlashAttention}: IO-aware tiling converts attention from memory-bound to compute-bound. 2--4$\times$ speedup.
  \item \textbf{Gradient checkpointing}: 10$\times$ less activation memory for 33\% more compute.
  \item \textbf{Profile before optimize}: Diagnose the bottleneck (compute, memory, or data) before choosing a technique.
  \item \textbf{Scaling trade-off}: Multi-GPU adds communication overhead. Exhaust single-machine optimizations first.
\end{itemize}

\end{frame}

\begin{frame}{References}
\note{[1 min] Point students to canonical papers.}

\small
\mlsysref{Micikevicius+18}{Micikevicius et al. ``Mixed Precision Training.'' ICLR 2018.}
\mlsysref{Dao+22}{Dao et al. ``FlashAttention: Fast and Memory-Efficient Exact Attention.'' NeurIPS 2022.}
\mlsysref{Chen+16}{Chen et al. ``Training Deep Nets with Sublinear Memory Cost.'' arXiv 2016.}
\mlsysref{Goyal+17}{Goyal et al. ``Accurate, Large Minibatch SGD.'' arXiv 2017.}
\mlsysref{Kingma+15}{Kingma \& Ba. ``Adam: A Method for Stochastic Optimization.'' ICLR 2015.}
\mlsysref{Brown+20}{Brown et al. ``Language Models are Few-Shot Learners.'' NeurIPS 2020.}

\end{frame}

\begin{frame}{Next Lecture: Data Selection}
\note{[1 min] Forward hook. Training taught us how to execute the optimization
loop efficiently. But we assumed the training data was given. The next
question: which data should we train on? Data selection strategies determine
what the model learns. Not all data is equally valuable.}

\footnotesize
\begin{columns}[c]
  \begin{column}{0.30\textwidth}
    \centering
    \begin{mlsyscard}{datastroke}
    {\large\bfseries Data Curation}\\[0.1cm]
    {\small Which examples matter most?}
    \end{mlsyscard}
  \end{column}
  \begin{column}{0.30\textwidth}
    \centering
    \begin{mlsyscard}{computestroke}
    {\large\bfseries Active Learning}\\[0.1cm]
    {\small Let the model choose its data.}
    \end{mlsyscard}
  \end{column}
  \begin{column}{0.30\textwidth}
    \centering
    \begin{mlsyscard}{routingstroke}
    {\large\bfseries Curriculum Design}\\[0.1cm]
    {\small Order matters for convergence.}
    \end{mlsyscard}
  \end{column}
\end{columns}

\vspace{0.3cm}
\centering
{\normalsize Training taught us \emph{how} to learn efficiently.}\\[0.1cm]
{\small\textbf{Data Selection asks: what should the model learn \emph{from}?}}

\end{frame}

\end{document}
